\textbf{¿Qué es?} En un árbol con raíz, el LCA de dos nodos \(u,v\) es el nodo más profundo que es ancestro de ambos a la vez (el punto donde "se cruzan" los caminos de \(u\) y de \(v\) hacia la raíz). \textit{¿Para qué sirve?} Distancia entre dos nodos: \(\text{dist}(u,v) = \text{profundidad}(u)+\text{profundidad}(v)-2\cdot\text{profundidad}(\text{LCA}(u,v))\); consultas sobre el camino entre dos nodos; verificar si un nodo es ancestro de otro. \textbf{Convención:} \texttt{grafo} = lista de adyacencia 0-indexada de un árbol \textbf{no dirigido} (\texttt{grafo[u]} = vecinos de \(u\)), \texttt{n} = número de nodos, \texttt{raiz} = nodo desde el que se cuelga el árbol (cualquiera, típicamente \texttt{0}).

\textit{¿Cuál versión usar?} \textbf{Binary Lifting}: más simple de programar bajo presión, \(O(\log n)\) por consulta, y de regalo te da "subir \(k\) niveles desde un nodo" (\(k\)-ésimo ancestro), útil en otros problemas de árboles. \textbf{Euler Tour + Sparse Table}: consulta \(O(1)\) (más rápida), pero más código — conviene si hay muchísimas consultas (\(\gtrsim 10^5\)). Ambas asumen que el árbol \textbf{no cambia} después de construir la estructura.
